\begin{lemma}
Let $0<\varepsilon\le1$ and let $k$ be a positive integer. There exists
a natural $A$ such that, for every natural $\gamma_0$, one can choose a
positive natural $\gamma\ge\gamma_0$ with the following property.
For every natural $D_{\rm nib}$ there exists a natural $D_0$, independent
of $\theta$, such that for every real $D\ge D_0$ and $0<\theta<1$
there are $s\in\mathbb N$ and natural sequences $N_i,K_i$ satisfying,
with $\ell_i=\lceil kN_i/\gamma\rceil$ and
$\rho=1-(1-\varepsilon/6)/\gamma$,
\[
D\le N_0,\qquad \ell_0\le K_0,
\]
and for every $i<s$,
\[
D_{\rm nib}\le N_i,\quad \varepsilon D/3\le N_i,\quad
K_i\le AN_i,\quad K_i\le K_{i+1},\quad
\rho N_i\le N_{i+1},\quad
(1-\theta-\varepsilon/6)\ell_i+K_{i+1}-K_i\ge\ell_{i+1}.
\]
Finally $K_s+kN_s+1\le(\theta+\varepsilon)kD$.
\end{lemma}
\begin{proof}
Set $A=\lceil6(1+\varepsilon)k/\varepsilon\rceil$. Choose a positive
natural $\gamma\ge\gamma_0$ with $1/\gamma\le\varepsilon/12$.
Then $0<\rho<1$. Given $D_{\rm nib}$, choose a natural $D_0$ so large
that every $D\ge D_0$ satisfies
\[
D\ge1,\quad D\ge60/\varepsilon,\quad
T:=\varepsilon D/3\ge D_{\rm nib},\quad (1-\rho)T\ge2,
\quad \frac{k\varepsilon^2D}{36\gamma}\ge
\frac{3k}{2}+3+\frac{k}{\gamma}.
\tag{1}
\]
Each condition is a lower bound on $D$ by a fixed real constant,
so the Archimedean property supplies this natural threshold; none
depends on $\theta$.
Put $B=\lceil D\rceil$, $N_0=B$. While $N_i\ge T$, set
$N_{i+1}=\lceil\rho N_i\rceil$. At such a step
$N_{i+1}<\rho N_i+1\le N_i-1<N_i$ by (1).
Thus there is a first index $s$ with $N_s<T$, since otherwise a natural
number would strictly decrease indefinitely. Also $B\ge D>T$, so $s>0$.
For $i\le s$ put
\[
K_i=\left\lceil(\theta+\varepsilon/2)k(B-N_i)+kB/\gamma\right\rceil.
\]
Extend both sequences constantly after $s$. Their defining ceiling
arguments are nonnegative, $N_i\le B$, $D\le B<D+1\le2D$,
and $K_0=\ell_0$. Monotonicity of the ceiling gives $K_i\le K_{i+1}$.
The degree and threshold inequalities in the statement follow immediately.

Here is the list-increment calculation, including its sign issue. Write
$c=1-\theta-\varepsilon/6>-\varepsilon/6$. If $c\ge0$, then
$c\ell_i\ge ckN_i/\gamma$. If $c<0$, use
$\ell_i\le kN_i/\gamma+1$ to obtain
$c\ell_i\ge ckN_i/\gamma-\varepsilon/6$.
The ceiling-difference bound and
$N_i-N_{i+1}\ge(1-\rho)N_i-1$ therefore give
\[
\begin{aligned}
c\ell_i+K_{i+1}-K_i
&\ge \frac{kN_i}{\gamma}
 [1+\varepsilon/2-(\varepsilon/6)(1+\theta+\varepsilon/2)]
 -(\theta+\varepsilon/2)k-1-\varepsilon/6\\
&\ge (1+\varepsilon/12)kN_i/\gamma
 -(\theta+\varepsilon/2)k-1-\varepsilon/6.
\end{aligned}
\]
In the last line $1+\theta+\varepsilon/2\le5/2$ was used.
Meanwhile $\ell_{i+1}\le k\rho N_i/\gamma+k/\gamma+1$.
Subtracting this upper bound from the preceding lower bound leaves at least
\[
\frac{k\varepsilon N_i}{12\gamma}-\frac{3k}{2}-2
-\frac{\varepsilon}{6}-\frac{k}{\gamma}\ge0,
\]
by $N_i\ge T$ and (1). This proves the list inequality even when $c<0$.
For $i<s$ we also have
\[
K_i\le(\theta+\varepsilon/2)kB+kB/\gamma+1
\le(1+7\varepsilon/12)kB+1
\le(1+\varepsilon)kB\le2(1+\varepsilon)kD\le AN_i.
\]
The third inequality follows from
$(5\varepsilon/12)kB\ge25$; the last uses $N_i\ge T$.
Finally, using $N_s<T$, $1/\gamma\le\varepsilon/12$, and
$0\le B-D<1$, we get
\[
\begin{aligned}
K_s+kN_s+1
&\le(\theta+\varepsilon/2)kB+kB/\gamma+2+k\varepsilon D/3\\
&\le(\theta+11\varepsilon/12)kD
 +(\theta+7\varepsilon/12)k(B-D)+2\\
&\le(\theta+11\varepsilon/12)kD+19k/12+2\\
&\le(\theta+\varepsilon)kD.
\end{aligned}
\]
Indeed $\varepsilon kD/12\ge5k\ge19k/12+2$ since $k\ge1$.
This includes the cost of rounding the initial real degree to $B$.
\end{proof}
